% Generated draft: requires mathematical and source review.
\documentclass[11pt]{article}
\usepackage[margin=25mm]{geometry}
\usepackage{fontspec}
\setmainfont{Latin Modern Roman}
\usepackage{amsmath,amssymb,mathtools,unicode-math}
\setmathfont{Latin Modern Math}
\usepackage{enumitem,needspace,xurl,hyperref}
\setlength{\parindent}{0pt}
\setlength{\parskip}{5pt}
\setlength{\emergencystretch}{3em}
\newcommand{\N}{\mathbb{N}}
\newcommand{\Z}{\mathbb{Z}}
\newcommand{\Q}{\mathbb{Q}}
\newcommand{\R}{\mathbb{R}}
\newcommand{\C}{\mathbb{C}}
\newcommand{\F}{\mathbb{F}}
\newcommand{\A}{\mathbb{A}}
\newcommand{\PP}{\mathbb{P}}
\newcommand{\E}{\mathbb{E}}
\newcommand{\eps}{\varepsilon}
\newcommand{\dd}{\,\mathrm{d}}
\begin{document}
\section*{UnsolvedMath ID: 30002665}
\textbf{Source problem code:} catalogue entry
\subsection{Definitions and mathematical statement}
For a precise provisional formulation, regard a convex \(n\)-gon \(Q\subset\R^2\) as a compact convex polygon with nonempty interior and exactly \(n\) vertices, including its boundary. These conventions are not explicitly specified in the supplied record and require confirmation.

An indexed family of translates of \(Q\) is a family \(\mathcal T=(Q+t_i)_{i\in I}\), where \(t_i\in\R^2\). Distinct indices may represent the same translate. For a positive integer \(m\), call \(\mathcal T\) an \(m\)-fold covering of the plane if
\[
\#\{i\in I:x\in Q+t_i\}\ge m
\qquad\text{for every }x\in\R^2.
\]
A decomposition into two coverings is a partition \(I=I_0\mathbin{\dot\cup}I_1\) such that
\[
\bigcup_{i\in I_0}(Q+t_i)=\R^2
\qquad\text{and}\qquad
\bigcup_{i\in I_1}(Q+t_i)=\R^2.
\]
Define
\[
m(Q)=\min\Bigl\{m\in\Z_{\ge1}:\text{every }m\text{-fold covering by translates of }Q
\text{ decomposes into two coverings}\Bigr\},
\]
with \(m(Q)=+\infty\) if this set is empty. The value \(+\infty\) is an explicit bookkeeping convention for possible nonexistence of the least multiplicity.

\textbf{Question.} Is it true that
\[
\forall n\in\Z,\ n>3,\quad
\exists M\in\Z_{\ge1}\quad
\forall\text{ convex }n\text{-gons }Q\subset\R^2,
\qquad m(Q)\le M?
\]
Thus the bound \(M=m(n)\) may depend on the number of vertices but must be independent of the shape and size of \(Q\).

\textbf{Unresolved scope qualification.} The displayed definition provisionally quantifies over arbitrary indexed families, allows repeated translates, and imposes no local finiteness restriction. The supplied record does not specify these choices, whether polygon boundaries are included, or whether degenerate polygons or redundant collinear vertices are allowed. Consequently this explicit interpretation requires source confirmation before being treated as the definitive intended formulation.
\subsection{Short English statement}
For each integer \(n>3\), is there a finite multiplicity depending only on \(n\) such that every covering of the plane at that multiplicity by translates of any convex \(n\)-gon can be partitioned into two coverings of the whole plane? The admissible covering families and polygon conventions require confirmation. This is a statement-only draft requiring human mathematical review.
\subsection{Sources}
\begin{itemize}
\item Local UnsolvedMath registry, permanent ID \(30002665\), source problem code \texttt{catalogue entry}, title: \emph{Uniform Cover Decomposition Bounds for Convex Polygons}. Collection: \emph{Oberwolfach Reports Open Problems}. The supplied proposed year is 2014; no individual proposer is supplied.
\item Supplied source citation: \emph{Discrete Geometry (2015)}, Oberwolfach Reports, pp.\ 2235--2296, DOI \texttt{10.4171/owr/2014/40}. \url{https://doi.org/10.4171/owr/2014/40}. The supplied background describes Oberwolfach Reports as documenting mathematical workshops and their open-problem discussions.
\item Supplied literature source, labeled ``Oberwolfach report'': \url{https://ems.press/content/serial-article-files/46530}.
\item Supplied literature source, labeled ``Pach--P\'{a}lv\"{o}lgyi, unsplittable coverings in the plane'': \url{https://doi.org/10.1016/j.aim.2016.07.011}.
\item Supplied literature assessment dated \texttt{2026-08-21}: the record labels the problem open and states that covering multiplicities are unbounded over all convex polygons, while the proposed bound for a fixed number of vertices remains unresolved even for quadrilaterals. This is supplied dated curation derived from the OWR material, not an independently verified literature assessment. Time-sensitive status and cited mathematical claims require independent verification.
\item The registry labels the statement \texttt{corrected\_verified}, explaining that only notation, typography, and TeX-delimiter cleanup was applied and that hypotheses, quantifiers, and subquestions were preserved. This supplied label does not constitute independent mathematical or source verification of this draft.
\end{itemize}
\end{document}
